
\setcounter{footnote}{0}
\SGARefTarget{sga3:xvi:chapter:xvi}\chapter*{Exposé XVI\\
Groups of Unipotent Rank Zero}
\addcontentsline{toc}{chapter}{Exposé XVI. Groups of Unipotent Rank Zero}
\markboth{EXPOSÉ XVI. GROUPS OF UNIPOTENT RANK ZERO}
  {EXPOSÉ XVI. GROUPS OF UNIPOTENT RANK ZERO}

\begin{center}
  \textit{by M. Raynaud}\textsuperscript{(*)}
\end{center}
\begingroup
\renewcommand{\thefootnote}{0}
\footnotetext{xy version of December 1, 2008.}
\renewcommand{\thefootnote}{*}
\footnotetext{Cf.\ the note on page 1 of Exposé XV.}
\endgroup

\SGARefTarget{sga3:xvi:section:1}\section*{1. An Immersion Criterion}
\addcontentsline{toc}{section}{1. An Immersion Criterion}

\SGARefTarget{sga3:xvi:subsection:1:1}%
\subsection*{1.1. Examples of Monomorphisms of Group Preschemes
That Are Not Immersions}
\addcontentsline{toc}{subsection}{\SGARefLink{sga3:xvi:subsection:1:1}{1.1}. Examples of Monomorphisms of
Group Preschemes That Are Not Immersions}

We shall construct a prescheme \(S\), two \(S\)-group preschemes
\(G\) and \(H\), and an \(S\)-group monomorphism
\[
  u:G\longrightarrow H
\]
that is not an immersion. The groups \(G\) and \(H\) will be of finite
presentation over \(S\), flat over \(S\), and \(G\) will even be
smooth over \(S\) (cf.\ Exposé VIII, the paragraph and note~\((*)\)
preceding \SGARefLink{sga3:viii:lemma:7:1}{7.1}; see also XVII, Appendix III, 4).

First take \(S\) to be the spectrum of a discrete valuation ring
\(A\) of unequal characteristic and residue characteristic \(2\).
Let \(t\) be the generic point of \(S\) and \(s\) the closed point.

\medskip
\noindent\textit{a)} Take for \(G\) the open subgroup of the constant
group
\[
  G'=(\mathbb Z/2\mathbb Z)_S
\]
obtained by removing the point of the closed fiber
\((\mathbb Z/2\mathbb Z)_s\) distinct from the identity element, and
take for \(H\) the group of multiplicative type
\((\mu_2)_S\) of square roots of unity (\SGARefLink{sga3:i:paragraph:4-4-4}{Exposé I, 4.4.4}). By the
choice of \(S\), \(H_t\) is isomorphic to
\((\mathbb Z/2\mathbb Z)_t\), whereas \(H_s\) is a radicial group.
There is an evident morphism \(G'\to H\), defined by the section
\((-1)\) of \(H\), hence a morphism \(u:G\to H\) which is a
monomorphism but not an immersion (otherwise it would necessarily be
an isomorphism; cf.\ \SGARefLink{sga3:viii:section:7}{Exposé VIII, 7}). Here \SGARefLink{sga3:viii:proposition:7:9}{Exposé VIII, 7.9} does not
apply, since \({}_{2}G=G\) is not finite over \(S\).

\medskip
\noindent\textit{b)} Let \(K\) denote the étale, nonseparated
\(S\)-group prescheme obtained from the identity group by
``doubling'' the unique point of the closed fiber. Take for \(H\) the
product group
\[
  (\mu_2)_S\times_S K.
\]
Let \(a\) (respectively \(b\)) be the unique section of \(\mu_2\)
(respectively \(K\)) over \(S\) distinct from the identity section,
and let \(h=(a,b)\) be the corresponding section of \(H\). The datum
of \(h\) defines an \(S\)-group morphism
\[
  u:G=(\mathbb Z/2\mathbb Z)_S\longrightarrow H.
\]
It is clear that \(u\) is a monomorphism and is not an immersion.
Here \SGARefLink{sga3:viii:proposition:7:9}{Exposé VIII, 7.9} does not apply, since \({}_{2}H=H\) is not
separated over \(S\).

\medskip
\noindent\SGAInlineRefTarget{sga3:xvi:subsection:1:1:item:c}\textit{c)} More interesting is the following example, in
which \(G\) is a smooth group with connected fibers (and even
\(G=(\Ga)_S\)); cf.\ \SGARefLink{sga3:viii:remark:7:10}{Exposé VIII, 7.10}.

Take for \(S\) the spectrum of a discrete valuation ring \(A\) of
equal characteristic \(2\), and let \(\pi\) be a uniformizer of
\(A\) and \(\mathfrak m\) the maximal ideal of \(A\).

Consider the additive group \((\Ga)_S\), the product group
\((\Ga\times\Ga)_S\) with ring \(A[X,Y]\), and let \(G_1\) be the
closed subgroup defined by
\begin{equation}
  X^2+X-\pi Y=0.
  \tag{1}\label{sga3:xvi:eq:1}
\end{equation}
The group \(G_1\) is smooth over \(S\); its generic fiber is
isomorphic to \((\Ga)_t\), and its closed fiber is isomorphic to the
product of \((\Ga)_s\) with \((\mathbb Z/2\mathbb Z)_s\). For the
\((\mathbb Z/2\mathbb Z)_s\)-factor one may take the subgroup \(N_s\)
of \((G_1)_s\) defined by
\begin{equation}
  Y=X^2+X=0.
  \tag{2}\label{sga3:xvi:eq:2}
\end{equation}

Let \(N'\) and \(N''\) be two group subschemes of \(G_1\), isomorphic
to \((\mathbb Z/2\mathbb Z)_S\), whose closed fibers coincide with
\(N_s\) and whose generic fibers are distinct. Thus \(N'\) and
\(N''\) are defined by two sections of \(G_1\) over \(S\), with
coordinates \((a',b')\) and \((a'',b'')\), respectively, such that
\begin{equation}
  \begin{cases}
    a'^2+a'-\pi b'=a''^2+a''-\pi b''=0,\\
    a'\text{ and }a''\equiv 1\pmod{\mathfrak m},
      \qquad a'\ne a'',\\
    b'\text{ and }b''\equiv 0\pmod{\mathfrak m}.
  \end{cases}
  \tag{3}\label{sga3:xvi:eq:3}
\end{equation}
(For example, one may take \((1,0)\) and
\((1+\pi^2,\pi+\pi^3)\).)

The quotient groups
\[
  G'=G_1/N'
  \qquad\text{and}\qquad
  G''=G_1/N''
\]
are representable (\SGARefLink{thm:sga3-V-4.1}{Exposé V, 4.1}). Let \(u\) be the canonical
morphism
\[
  u:G_1\longrightarrow G'\times_S G'',
\]
and let \(H\) be the group subscheme of \(G'\times_S G''\) equal to
the schematic closure in \(G'\times_S G''\) of
\(u_t((G_1)_t)\), so that \(u\) factors through \(H\). The kernel of
\(u\) is
\[
  K=N'\cap N'',
\]
whose generic fiber is the identity group and whose closed fiber is
equal to \(N_s\); in particular, \(K\) is not flat over \(S\). On the
generic fiber, \(u_t\) is therefore an isomorphism from
\((G_1)_t\) onto \(H_t\). On the other hand, \(H_s\) is not smooth.
Indeed, if \(H_s\) were smooth, then, since \(u_s\) is a morphism with
finite kernel and \((G_1)_s\) and \(H_s\) have the same dimension,
namely \(1\), \(u_s\) would be flat. Since \(G_1\) and \(H\) are flat,
\(u\) would then be flat, and consequently \(\Ker u\) would be flat
over \(S\), which it is not. It is clear that the restriction of
\(u\) to the identity component \(G\) of \(G_1\) is a monomorphism:
the fibers of \(K\cap G\) are identity groups, so \(K\cap G\) is the
identity group (Exposé VI\textsubscript{B}, \SGARefLink{lem:sga3-VIB-2.9}{2.9}). But it is not an
immersion, for otherwise it would be an open immersion and \(H_s\)
would be smooth. Here \SGARefLink{sga3:viii:proposition:7:9}{Exposé VIII, 7.9} does not apply, since
\({}_{2}G=G\) is not finite over \(S\).

For lovers of equations, in the preceding example one may take for
\(H\) the closed subgroup of \((\Ga\times_S\Ga)_S\) with ring
\(A[V,W]/(F)\), where
\[
  F=(a''b'-b''a')(a''^2V-a'^2W)-(a''V-a'W)^2
\]
(where \(a',a'',b',b''\) satisfy \eqref{sga3:xvi:eq:3}). For \(G\)
one takes the group \((\Ga)_S\) with ring \(A[T]\), and for the
morphism \(u:G\to H\) the \(S\)-morphism defined by
\[
  V\longmapsto a'(a'T+\pi T^2),
  \qquad
  W\longmapsto a''(a''T+\pi T^2).
\]

\medskip
\noindent\SGARefTarget{sga3:xvi:remark:1:2}%
\textbf{Remark 1.2.} --- The preceding construction is inspired by
the method of Koizumi--Shimura.\footnote{Editor's note:
S.~Koizumi and G.~Shimura, ``Specialization of abelian varieties,''
\emph{Sci.\ Papers Coll.\ Gen.\ Ed.\ Univ.\ Tokyo} 9 (1959),
187--211.} It does not apply when \(S\) is the spectrum of a ring of
unequal characteristic, because the points of finite order of \(G\)
are then ``too rigid.''

\SGARefTarget{sga3:xvi:subsection:1:3}%
\subsection*{1.3. Statement of the Immersion Criterion}
\addcontentsline{toc}{subsection}{1.3. Statement of the Immersion Criterion}

\medskip
\noindent\SGARefTarget{sga3:xvi:theorem:1:3}%
\textbf{Theorem 1.3.} --- \textit{Let \(S\) be a locally noetherian
prescheme, let \(G\) be an \(S\)-group prescheme, smooth and of finite
type over \(S\), possessing Cartan subgroups locally for the fpqc
topology (\SGARefLink{sga3:xv:definition:6-1}{Exposé XV, 6.1} and \SGARefLink{sga3:xv:theorem:7-3}{7.3}~(i)), let \(H\) be an \(S\)-group
prescheme locally of finite type, and let \(u:G\to H\) be an
\(S\)-group monomorphism. Then:}

\begin{enumerate}
\item[\textit{a)}] \textit{If \(G\) has connected fibers, then \(u\)
is an immersion if and only if, for every \(S\)-scheme \(S'\) that is
the spectrum of a complete discrete valuation ring with algebraically
closed residue field, and every Cartan subgroup \(C\) of \(G_{S'}\),
the restriction of \(u_{S'}\) to \(C\) is an immersion.}

\item[\textit{b)}] \textit{The morphism \(u\) is an immersion
(respectively a closed immersion) if and only if, for every \(S'\) as
above and every Cartan subgroup \(C\) of the identity component
\((G_{S'})^0\) of \(G_{S'}\), the restriction of \(u_{S'}\) to the
normalizer \(N\) of \(C\) in \(G_{S'}\) (\SGARefLink{sga3:xi:corollary:6:11}{Exposé XI, 6.11}) is an
immersion (respectively a closed immersion).}
\end{enumerate}

Before proving \SGARefLink{sga3:xvi:theorem:1:3}{1.3}, let us state some applications.

\medskip
\noindent\SGARefTarget{sga3:xvi:corollary:1:4}%
\textbf{Corollary 1.4.} --- \textit{Let \(S\) be a prescheme, let
\(G\) be an \(S\)-group prescheme, smooth and of finite presentation
over \(S\), of unipotent rank zero (\SGARefLink{sga3:xv:definition:6:1:ter}{Exposé XV, 6.1 ter}), possessing
maximal tori locally for the fpqc topology, let \(H\) be an
\(S\)-group prescheme, and let \(u:G\to H\) be an \(S\)-group
monomorphism. Suppose in addition either that \(H\) is of finite
presentation over \(S\), or that \(S\) is locally noetherian and
\(H\) is locally of finite type. Then:}

\begin{enumerate}
\item[\textit{a)}] \textit{If \(G\) has connected fibers, \(u\) is
an immersion.}

\item[\textit{b)}] \textit{If \(H\) is separated over \(S\), and if
for every \(S'\) over \(S\) and every maximal torus \(T\) of
\(G_{S'}\), the Weyl group}
\[
  W=
  \Norm_{G_{S'}}(T)\big/\Centr_{G_{S'}^0}(T)
\]
\textit{is representable (a condition always satisfied when \(G\)
has connected fibers; \SGARefLink{sga3:xv:proposition:7-1}{Exposé XV, 7.1}~(iv)) and finite over \(S\), then
\(u\) is a closed immersion.}
\end{enumerate}

\medskip
\noindent\SGARefTarget{sga3:xvi:corollary:1:5}%
\textbf{Corollary 1.5.} --- \textit{Let \(S\) be a prescheme, let
\(G\) be an \(S\)-group prescheme, smooth and of finite presentation
over \(S\), with connected fibers, let \(H\) be an \(S\)-group
prescheme, and let \(u:G\to H\) be an \(S\)-group monomorphism.
Suppose either that \(H\) is of finite presentation over \(S\), or
that \(S\) is locally noetherian and \(H\) is locally of finite type.
Then:}

\begin{enumerate}
\item[\textit{a)}] \textit{If \(G\) is reductive, that is, affine
over \(S\) with reductive fibers (cf.\ Exposé XIX), \(u\) is an
immersion, and a closed immersion if \(H\) is separated.}

\item[\textit{b)}] \textit{If \(G\) has affine, solvable, connected
fibers of unipotent rank zero, \(u\) is an immersion, and a closed
immersion if \(H\) is separated.}
\end{enumerate}

\medskip
\noindent\textit{Proof of \SGARefLink{sga3:xvi:corollary:1:4}{1.4} from \SGARefLink{sga3:xvi:theorem:1:3}{1.3}.}

\smallskip
\noindent\textit{Reduction to the noetherian case.}
If \(S\) is locally noetherian, the reduction is immediate, since the
properties to be proved are local on \(S\). In the second case \(G\)
and \(H\) are of finite presentation over \(S\). After shrinking
\(S\), we may suppose that \(S\) is affine. By Exposé
VI\textsubscript{B}, \S10, there exist a noetherian scheme \(S_0\), an
\(S_0\)-group prescheme \(G_0\), smooth over \(S_0\) (with connected
fibers in case~a)) and of finite type over \(S_0\), an
\(S_0\)-group prescheme \(H_0\) of finite type (separated in case~b)),
and an \(S_0\)-group morphism
\[
  u_0:G_0\longrightarrow H_0
\]
such that \(G,H,u\) are obtained from \(G_0,H_0,u_0\) by the base
extension \(S\to S_0\). The assertion that \(u\) is a monomorphism
means that \(\Ker u\) is the identity group; we may therefore suppose
that \(u_0\) is a monomorphism.

Since \(G\) has unipotent rank \(\rho_u=0\) and possesses maximal
tori locally for the fpqc topology, the abelian rank \(\rho_{ab}\) of
the fibers of \(G\) is locally constant (\SGARefLink{sga3:xv:corollary:8-18}{Exposé XV, 8.18}). But
\(\rho_u\) and \(\rho_{ab}\) are locally constructible functions
(Exposé XV, 6.3 bis). A standard argument (cf.\ EGA IV, 8.3.4) shows
that one may choose \(S_0\) and \(G_0\) so that the unipotent rank
(respectively the abelian rank) of the fibers of \(G_0\) is zero
(respectively locally constant). Then \(G_0\) possesses maximal tori
locally for the fpqc topology (\SGARefLink{sga3:xv:corollary:8-18}{Exposé XV, 8.18}), and the functor
\(\mathcal{TM}_{G_0}\) of maximal tori of \(G_0\) is representable by
an \(S_0\)-scheme \(X_0\) of finite type over \(S_0\) (\SGARefLink{sga3:xv:theorem:8-15}{Exposé XV,
8.15}). Let \(T_0\) be the ``universal'' maximal torus of
\((G_0)_{X_0}\), let \(X=X_0\times_{S_0}S\), and let
\(T=T_0\times_{S_0}S\), the universal maximal torus for \(G\). By
hypothesis, in case~b), the Weyl group relative to \(T\) is
representable and finite over \(X\). These two properties are
compatible with projective limits of preschemes (Exposé
VI\textsubscript{B}, 10.1~iii, and EGA IV, 8.10.5). We may therefore
choose \(S_0\) so that the Weyl group of \(T_0\) is finite over
\(X_0\). Under these conditions, to prove \SGARefLink{sga3:xvi:corollary:1:4}{1.4} we may replace
\(S,G,u,H\) by \(S_0,G_0,u_0,H_0\), and hence suppose that \(S\) is
noetherian.

Use the valuative criterion furnished by \SGARefLink{sga3:xvi:theorem:1:3}{1.3}. We are reduced to the
case where \(S\) is the spectrum of a discrete valuation ring and
where \(G^0\) possesses a Cartan subgroup \(C\).

\smallskip
\noindent\textit{Proof of \SGARefLink{sga3:xvi:corollary:1:4}{1.4} a).}
We must show that the restriction of \(u\) to \(C\) is an immersion
(\SGARefLink{sga3:xvi:theorem:1:3}{1.3}~a)), which reduces us to the case where \(G=C\) has connected,
nilpotent fibers. Standard reductions (cf.\ the proof of \SGARefLink{sga3:viii:lemma:7:1}{Exposé VIII,
7.1}) reduce us to the case where \(H\) is flat over \(S\), \(u\) is
an isomorphism on the generic fiber, and \(u\) is an isomorphism of
the underlying spaces on the closed fiber. The group \(H\) then has
connected fibers and is therefore separated over \(S\) (Exposé
VI\textsubscript{B}, \SGARefLink{prop:sga3-VIB-5.2}{5.2}). Grant for a moment the following lemma.

\medskip
\noindent\SGARefTarget{sga3:xvi:lemma:1:6}%
\textbf{Lemma 1.6.} --- \textit{Let \(S\) be a prescheme and let
\(G\) be an \(S\)-group prescheme, smooth over \(S\), with connected,
nilpotent fibers and unipotent rank zero. Then:}
\begin{enumerate}
\item[\textit{i)}] \textit{\(G\) is commutative.}
\item[\textit{ii)}] \textit{For every integer \(n>0\),
\({}_{n}G\) is an \(S\)-group prescheme, flat and quasi-finite over
\(S\), and it is finite over \(S\) if and only if the abelian rank,
or the reductive rank, of \(G\) is locally constant on \(S\).}
\item[\textit{iii)}] \textit{For every integer \(q>0\) invertible on
\(S\), the family of subgroups \({}_{q^n}G\) is universally
schematically dense in \(G\) relative to \(S\) (EGA IV, 11.10.8).}
\end{enumerate}

\SGARefLink{sga3:xvi:lemma:1:6}{Lemma 1.6} applies to \(G\), and since \(G\) possesses maximal tori
locally, the reductive rank of the fibers of \(G\) is locally
constant. Thus, by \SGARefLink{sga3:xvi:lemma:1:6}{1.6}~ii), \({}_{n}G\) is finite over \(S\) for
every integer \(n>0\). It follows from \SGARefLink{sga3:viii:proposition:7:9}{Exposé VIII, 7.9} that \(u\) is
an immersion.

\medskip
\noindent\textit{Proof of \SGARefLink{sga3:xvi:lemma:1:6}{1.6}.}
Let us first examine the case where \(S\) is the spectrum of a field.

\medskip
\noindent\SGARefTarget{sga3:xvi:lemma:1:7}%
\textbf{Lemma 1.7.} --- \textit{Let \(k\) be a field of
characteristic \(p\), and let \(G\) be a smooth, connected, nilpotent
algebraic \(k\)-group of unipotent rank zero. Then \(G\) is a
commutative group, an extension of an abelian variety \(A\) by a
torus \(T\). For every integer \(n>0\), \({}_{n}G\) is a finite
group, étale if \((n,p)=1\), defined by a \(k\)-algebra of rank
\(n^{\rho_r+2\rho_{ab}}\). If \((q,p)=1\), the family of subgroups
\({}_{q^n}G\), \(n\in\mathbb N\), is schematically dense in \(G\).}

\smallskip
\noindent\textit{Proof of \SGARefLink{sga3:xvi:lemma:1:7}{1.7}.}
Let \(Z\) be the center of \(G\). The group \(G/Z\) is affine
(\SGARefLink{sga3:xii:6.1}{Exposé XII, 6.1}), smooth, connected, and of unipotent rank zero;
hence it is a torus (BIBLE 4, th.~4). It follows that \(G\) is
commutative (Exposé XII, 6.4). If \(T\) is the unique maximal torus
of \(G\) (Exposé XV, 3.4), Chevalley's theorem (Séminaire Bourbaki
1956/57, no.~145) immediately shows that \(G\) is an extension of an
abelian variety \(A\) by \(T\). For every integer \(n>0\), raising to
the \(n\)-th power is an epimorphism on \(T\), whence the exact
sequence
\[
  0\longrightarrow {}_{n}T\longrightarrow {}_{n}G
  \longrightarrow {}_{n}A\longrightarrow 0.
\]
A classical theorem of Weil (A.~Weil, \emph{Variétés abéliennes et
courbes algébriques}, \S IX, th.~33, cor.~1) says that \({}_{n}A\)
is a finite group defined by a \(k\)-algebra of rank
\(n^{2\rho_{ab}}\). Since \({}_{n}T\) is a finite group of rank
\(n^{\rho_r}\), the asserted structure of \({}_{n}G\) follows.

Now let \(H\) be the smallest closed subscheme of \(G\) containing
\({}_{q^n}G\) for every \(n\). It follows from the preceding
discussion and \SGARefLink{sga3:xv:lemma:4:6}{Exposé XV, 4.6} that \(H\) is a smooth, connected
subgroup, hence of the same type as \(G\). Raising to the \(q\)-th
power on \(H\) is an epimorphism (since \({}_{q}H\) is finite), so
there is an exact sequence
\[
  0\longrightarrow {}_{q}H\longrightarrow {}_{q}G
  \longrightarrow {}_{q}(G/H)\longrightarrow 0.
\]
It follows that \({}_{q}(G/H)=0\), and therefore \(G/H=0\). Thus the
subgroups \({}_{q^n}G\) are schematically dense in \(G\).

\smallskip
\noindent\textit{Continuation of the proof of \SGARefLink{sga3:xvi:lemma:1:6}{1.6}~i).}
To show that \(G\) is commutative, the usual procedure reduces us
first to the case where \(S\) is noetherian, and then to the case
where \(S\) is the spectrum of a local ring with closed point \(s\).
The center of \(G\) is represented by a closed group subscheme \(Z\)
of \(G\) (\SGARefLink{sga3:xi:corollary:6:11}{Exposé XI, 6.11}). To prove that \(Z=G\), it suffices to
prove this after reduction modulo every power of the maximal ideal of
the ring of \(S\): \(Z\) will then be an open subgroup of \(G\), and
hence equal to \(G\), since \(G\) has connected fibers. This reduces
us to the case where \(S\) is local Artinian.

Let \(q\) be an integer invertible on \(S\). For every integer \(n\),
\({}_{q^n}G_s\) is an étale, central subgroup of multiplicative type
of \(G_s\) (\SGARefLink{sga3:xvi:lemma:1:7}{1.7}). Since \(G\) is smooth,
\({}_{q^n}(G_s)\) lifts to an étale, central \(S\)-group subscheme
\(M_n\) of \(G\) (\SGARefLink{sga3:xv:proposition:1:2}{Exposé XV, 1.2}). The family of flat subgroups
\(M_n\), \(n\in\mathbb N\), is then schematically dense in \(G\)
(1.7 and EGA IV, 11.10.9). Since \(Z\) is closed in \(G\) and
contains every \(M_n\), we have \(Z=G\).

\smallskip
\noindent\textit{Proof of \SGARefLink{sga3:xvi:lemma:1:6}{1.6}~ii).}
To see that \({}_{n}G\) is flat and quasi-finite over \(S\), it
suffices to show that the \(n\)-th power map on \(G\) is flat and
quasi-finite. Since \(G\) is flat over \(S\), we are reduced to the
case where \(S\) is the spectrum of a field (EGA IV, 11.3.10), and
we observed in the proof of 1.7 that the \(n\)-th power map is an
epimorphism. Moreover, \({}_{n}G\) is separated over \(S\), since
\(G\) is separated over \(S\) (Exposé VI\textsubscript{B}, \SGARefLink{prop:sga3-VIB-5.2}{5.2}).
For \({}_{n}G\) to be finite over \(S\), it is necessary and
sufficient that the fibers of \({}_{n}G\) be spectra of finite
algebras of locally constant rank. This is seen immediately by
reducing to the case where \(S\) is the spectrum of a henselian local
ring. By \SGARefLink{sga3:xvi:lemma:1:7}{1.7}, this condition is equivalent to saying that the abelian
rank, or the reductive rank, of the fibers of \(G\) is locally
constant.

\smallskip
\noindent\textit{Proof of \SGARefLink{sga3:xvi:lemma:1:6}{1.6}~iii).}
To see that the family of subgroups
\({}_{q^n}G\), \(n\in\mathbb N\), is universally schematically dense
in \(G\), we again reduce to the noetherian case. In view of 1.7 and
1.6~ii), it suffices to apply EGA IV, 11.10.9.

\medskip
\noindent\textit{Proof of \SGARefLink{sga3:xvi:corollary:1:4}{1.4} b).}
We have reduced to the case where \(S\) is the spectrum of a discrete
valuation ring and where \(G^0\) possesses a Cartan subgroup \(C\).
Let
\[
  N=\Norm_G(C)=\Norm_G(T),
\]
where \(T\) is the unique maximal torus of \(C\) (\SGARefLink{sga3:xii:7.1}{Exposé XII,
7.1}~a) and~b)). Since \(H\) is separated over \(S\), it follows from
\SGARefLink{sga3:xvi:lemma:1:6}{1.6}~ii) and \SGARefLink{sga3:viii:proposition:7:12}{Exposé VIII, 7.12} that the restriction of \(u\) to \(C\)
is a closed immersion. On the other hand, to prove that \(u\) is a
closed immersion, it suffices to prove this for \(u|_N\) (\SGARefLink{sga3:xvi:theorem:1:3}{1.3}~b)).
Since, by hypothesis, \(W=N/C\) is represented by a finite group
scheme, the assertion follows from the next lemma, applied to the
exact sequence
\[
  0\longrightarrow C\longrightarrow N\longrightarrow W
  \longrightarrow 1.
\]
(Recall that a proper immersion is a closed immersion.)

\medskip
\noindent\SGARefTarget{sga3:xvi:lemma:1:8}%
\textbf{Lemma 1.8.} --- \textit{Let \(S\) be a prescheme and let
\(G\) be an \(S\)-group prescheme of finite presentation over \(S\),
an extension of an \(S\)-group prescheme \(G''\), proper and of finite
presentation over \(S\), by an \(S\)-group prescheme \(G'\), flat and
of finite presentation over \(S\):}
\[
  1\longrightarrow G'\longrightarrow G\longrightarrow G''
  \longrightarrow 1.
\]
\textit{Let \(H\) be an \(S\)-group prescheme of finite presentation
over \(S\) (or locally of finite presentation if \(S\) is locally
noetherian), and let \(u:G\to H\) be an \(S\)-group morphism. If the
restriction of \(u\) to \(G'\) is proper, then \(u\) is proper.}

\smallskip
\noindent\textit{Proof of \SGARefLink{sga3:xvi:lemma:1:8}{1.8}.}
As usual, we reduce to the case where \(S\) is noetherian (Exposé
VI\textsubscript{B}, \S10, and Exposé XV, 6.2), and we may then apply
the valuative criterion of properness (EGA II, 7.3.8). We therefore
suppose that \(S\) is the spectrum of a discrete valuation ring
\(A\), with closed point \(s\) and generic point \(t\). Let
\(\bar x\in G(t)\) and \(h\in H(S)\) satisfy
\[
  u_t(\bar x)=h(t).
\]
We must show that \(\bar x\) comes from a unique element
\(x\in G(S)\). It is enough to prove the existence and uniqueness of
\(x\) after a faithfully flat extension of the discrete valuation
ring \(A\). Let \(\bar y\) be the image of \(\bar x\) in \(G''(t)\).
Since \(G''\) is proper over \(S\), \(\bar y\) comes from a unique
element \(y\in G''(S)\). Let \(X\) be the inverse image of \(y\) in
\(G\). The prescheme \(X\) is faithfully flat over \(S\), since
\(G'\) is faithfully flat over \(S\), as is the morphism
\(G\to G''\) (Exposé VI\textsubscript{B}, \S9). After a faithfully
flat change of discrete valuation ring, we may suppose that \(X\)
has a section \(\ell\) over \(S\) (EGA IV, 14.5.8). Replacing
\(\bar x\) by \(\bar x\ell_t^{-1}\), we may suppose that
\(\bar x\in G'(t)\). The existence and uniqueness of \(x\) now
follow from the properness of \(u|_{G'}\).

\medskip
\noindent\textit{Proof of \SGARefLink{sga3:xvi:corollary:1:5}{1.5}.}

\smallskip
\noindent\textit{a)}
We shall see in Exposé XIX that, if \(G\) is reductive, then \(G\)
possesses maximal tori locally for the faithfully flat topology, has
a finite Weyl group, and has unipotent rank zero. Assertion~a)
therefore follows from \SGARefLink{sga3:xvi:corollary:1:4}{1.4}.

\smallskip
\noindent\textit{b)}
If \(G\) has affine fibers and unipotent rank zero, then \(G\)
possesses maximal tori locally for the fpqc topology (\SGARefLink{sga3:xv:corollary:8-18}{Exposé XV,
8.18}). It remains only to apply \SGARefLink{sga3:xvi:corollary:1:4}{1.4}, observing that, since \(G\) has
connected, smooth, affine, solvable fibers, its Weyl group is the
identity group (BIBLE 6, th.~1, cor.~3).

\medskip
\noindent\textit{Proof of \SGARefLink{sga3:xvi:theorem:1:3}{1.3} a).}
Since \(u\) is already a monomorphism, to prove that \(u\) is an
immersion it suffices to prove that \(u\) is proper at the points of
\(u(G)\) (EGA IV, 15.7.1); for this we may use the valuative
criterion of local properness (EGA IV, 15.7.5). We are therefore
reduced to the case where \(S\) is the spectrum of a complete
discrete valuation ring with algebraically closed residue field,
closed point \(s\), and generic point \(t\). Since \(G\) is flat over
\(S\), standard reductions (cf.\ the proof of \SGARefLink{sga3:viii:lemma:7:1}{Exposé VIII, 7.1})
reduce us to the case where \(H\) is flat over \(S\), \(u_t\) is an
isomorphism, and \(u_s\) is an isomorphism of underlying spaces. To
say that \(u\) is an immersion is then equivalent to saying that
\(u\) is an isomorphism.

By hypothesis, \(G\) possesses Cartan subgroups locally for the fpqc
topology, so we may speak of the open subset \(U\) of regular points
of \(G\) (\SGARefLink{sga3:xv:theorem:7-3}{Exposé XV, 7.3}~i) and~ii)).

\medskip
\noindent\SGARefTarget{sga3:xvi:lemma:1:9}%
\textbf{Lemma 1.9.} --- \textit{Under the preceding hypotheses, for
\(u\) to be an immersion it suffices that \(u|_U\) be an immersion.}

Indeed, to say that \(u|_U\) is an immersion means that there exist
an open subscheme \(V\) of \(H\) and a closed subscheme \(F\) of
\(V\) such that \(u|_U\) factors through \(F\) and induces an
isomorphism from \(U\) onto \(F\). Since \(H_t\), and therefore
\(V_t\), is irreducible, one necessarily has \(V_t=F_t\). Then \(F\)
contains the schematic closure of \(V_t\) in \(V\), which equals
\(V\) because \(H\) is flat over \(S\). Thus \(V=F\). It follows
that \(u|_U\) is an open immersion, and hence that \(u\) is an open
immersion (Exposé VI\textsubscript{B}, \SGARefLink{cor:sga3-VIB-2.6}{2.6}).

It remains to prove that \(u|_U\) is an immersion. For this we apply
the valuative criterion of local properness. After replacing \(S\)
by the spectrum of a discrete valuation ring faithfully flat over
\(S\), we must show the following: if \(h\) is a section of \(H\)
over \(S\) whose image \(h(S)\) is contained in \(u(U)\), then \(h\)
is the image of a section \(g\) of \(U\) over \(S\). It suffices to
show that \(h\) is contained in the image of a Cartan subgroup \(C\)
of \(G\). Indeed, by hypothesis \(u|_C\) is an immersion, so \(h\)
is the image of a section \(g\) of \(C\), necessarily a section of
\(U\).

Let \(a\in U(s)\) and \(b\in U(t)\) satisfy, respectively,
\[
  u_s(a)=h(s),
  \qquad
  u_t(b)=h(t).
\]
Since \(u_t\) is an isomorphism, \(h(t)\) is a regular point of
\(H_t\), and is therefore contained in a unique Cartan subgroup
\(D_t\) of \(H_t\) (\SGARefLink{rem:sga3-VIB-3.2}{Exposé XIII, 3.2}). Let \(D\) be the schematic
closure of \(D_t\) in \(H\).

\medskip
\noindent\SGARefTarget{sga3:xvi:lemma:1:10}%
\textbf{Lemma 1.10.} ---
\begin{enumerate}
\item[\textit{i)}] \textit{\(D_s\) is a nilpotent algebraic group
(Exposé VI\textsubscript{B}, \S8).}
\item[\textit{ii)}] \textit{\(\dim D_s=\nu\), where \(\nu\) is the
nilpotent rank of \(G\), and also the nilpotent rank of
\((H_s)_{\mathrm{red}}\).}
\item[\textit{iii)}] \textit{\(h(s)\in D_s(\kappa(s))\).}
\end{enumerate}

\noindent\textit{i)}
Since \(H\) is separated (Exposé VI\textsubscript{B}, \SGARefLink{prop:sga3-VIB-5.2}{5.2}), so is
\(D\). Moreover, \(D\) is flat over \(S\) and its generic fiber is
nilpotent. It follows from Exposé VI\textsubscript{B}, \SGARefLink{prop:sga3-VIB-8.4}{8.4} that
\(D_s\) is nilpotent.

\smallskip
\noindent\textit{ii)}
Since \((H_s)_{\mathrm{red}}\simeq G_s\), these two groups have the
same nilpotent rank. Since \(G\) possesses Cartan subgroups locally
for the fpqc topology, \(G_s\) and \(G_t\) have the same nilpotent
rank \(\nu\). Finally, since \(D\) is flat over \(S\),
\[
  \dim D_s=\dim D_t=\nu
\]
(Exposé VI\textsubscript{B}, \S4).

\smallskip
\noindent\textit{iii)}
Since \(h_t\) factors through \(D_t\), the section \(h\) plainly
factors through \(D\), which proves iii).

The next lemma now shows that \((D_s)_{\mathrm{red}}\) is a Cartan
subgroup of
\[
  (H_s)_{\mathrm{red}}\simeq G_s.
\]

\medskip
\noindent\SGARefTarget{sga3:xvi:lemma:1:11}%
\textbf{Lemma 1.11.} --- \textit{Let \(G\) be a smooth, connected
algebraic group over an algebraically closed field \(k\), and let
\(D\) be a smooth, nilpotent algebraic subgroup containing a regular
element \(a\in G(k)\). Then \(D^0\) is contained in a Cartan subgroup
of \(G\). If, in addition, \(\dim D\) equals the nilpotent rank of
\(G\), then \(D\) is a Cartan subgroup of \(G\), and in particular is
connected.}

Let \(Z\) be the center of \(G\), let \(G'=G/Z\), which is affine
(\SGARefLink{sga3:xii:6.1}{Exposé XII, 6.1}), and let \(a'\) and \(D'\) be the images of \(a\)
and \(D\) in \(G'\). The correspondence between Cartan subgroups of
\(G\) and Cartan subgroups of \(G'\) (\SGARefLink{sga3:xii:6.6}{Exposé XII, 6.6}~e)) shows at
once that \(a'\) is a regular element of \(G'\) and that it suffices
to prove the lemma for the pair \((D',G')\). Thus we may suppose that
\(G\) is affine.

Let \(s\) be the semisimple component of \(a\). It belongs to
\(D(k)\) (BIBLE 4, th.~3) and is regular (BIBLE 7, th.~2, cor.~1).
By BIBLE 6, th.~2, \(s\) centralizes \(D^0\); hence \(D^0\) is
contained in the connected centralizer of \(s\), which is a Cartan
subgroup of \(G\) (BIBLE 7, th.~2). If now
\(\dim D\) equals the nilpotent rank of \(G\), then \(D^0\) is a
Cartan subgroup of \(G\), equal to \(\Centr_G(T)\), where \(T\) is
the unique maximal torus of \(D^0\) (\SGARefLink{sga3:xii:6.6}{Exposé XII, 6.6}). But \(D\) is
nilpotent, and therefore centralizes \(T\) (BIBLE 6, th.~2), so
\(D=D^0\).

Return now to the group \(G\). The group
\[
  C_t=u_t^{-1}(D_t)
\]
is the unique Cartan subgroup of \(G_t\) containing \(b\). Let \(C\)
be the schematic closure of \(C_t\) in \(G\).

By functoriality of schematic closure, \(u|_C\) factors through
\(D\), and therefore \(u_s(C_s)\subset D_s\). Since \(u_s\) is an
isomorphism from \(G_s\) onto \((H_s)_{\mathrm{red}}\), since
\[
  \dim C_s=\dim C_t=\nu=\dim D_s
\]
by \SGARefLink{sga3:xvi:lemma:1:10}{1.10}, and since \(D_s\) is connected by \SGARefLink{sga3:xvi:lemma:1:11}{1.11}, \(u_s\) induces an
isomorphism from \((C_s)_{\mathrm{red}}\) onto
\((D_s)_{\mathrm{red}}\). Thus \((C_s)_{\mathrm{red}}\) is a Cartan
subgroup of \(G_s\). The next lemma shows that \(C\) itself is a
Cartan subgroup of \(G\).

\medskip
\noindent\SGARefTarget{sga3:xvi:lemma:1:12}%
\textbf{Lemma 1.12.} --- \textit{Let \(S\) be the spectrum of a
discrete valuation ring, let \(G\) be a smooth \(S\)-group prescheme
of finite type, and let \(C\) be a group subscheme of \(G\), flat
over \(S\), such that the generic fiber \(C_t\) is a Cartan subgroup
of \(G_t\) and the geometric reduced closed fiber
\((C_{\bar s})_{\mathrm{red}}\) is a Cartan subgroup of
\(G_{\bar s}\). Then \(C\) is a Cartan subgroup of \(G\).}

We must prove that \(C\) is smooth over \(S\), and it suffices to
establish this after a faithfully flat extension of the base. The
hypotheses on \(C\) imply that the nilpotent rank of the fibers of
\(G\) is constant; hence we may speak of the open subset \(U\) of
regular points of \(G\) (\SGARefLink{sga3:xv:theorem:7-3}{Exposé XV, 7.3}). Since
\((C_{\bar s})_{\mathrm{red}}\) is a Cartan subgroup of
\(G_{\bar s}\), \(C_{\bar s}\) contains a regular point
\(a_{\bar s}\) of \(G_{\bar s}\). Since \(C\) is flat over \(S\),
after a faithfully flat extension of the base if necessary
(EGA IV, 14.5.8), we may suppose that \(a_{\bar s}\) is an element
of \(G(s)\) which lifts to a section \(a\in C(S)\). Since \(U\) is
open and contains \(a(s)\), it contains \(a\). Let \(C'\) be the
unique Cartan subgroup of \(G\) containing \(a\) (\SGARefLink{sga3:xv:theorem:7-3}{Exposé XV, 7.3}).
Necessarily \(C_t=C'_t\), and consequently \(C'\) and \(C\) coincide
with the identity component of the schematic closure of \(C_t\) in
\(G\) (observe that both \(C'\) and \(C\) are flat over \(S\)).

\smallskip
\noindent\textit{End of the proof of \SGARefLink{sga3:xvi:theorem:1:3}{1.3} a).}
By hypothesis, the restriction of \(u\) to the Cartan subgroup \(C\)
of \(G\) is an immersion. It is clear that \(h(S)\) is contained in
\(u(C)=D\), and hence \(h\) is the image of a section \(g\in C(S)\).
\hfill\(\square\)

\medskip
\noindent\textit{Proof of \SGARefLink{sga3:xvi:theorem:1:3}{1.3} b).}
We shall again use the valuative criterion of local properness
(EGA IV, 15.5) in the case of an immersion (respectively, the
valuative criterion of properness (EGA II, 7.3.6) in the case of a
closed immersion). This reduces us to the case where \(S\) is the
spectrum of a complete discrete valuation ring with algebraically
closed residue field. Let \(s\) be the closed point and \(t\) the
generic point of \(S\). Replacing \(H\) by the schematic closure in
\(H\) of \(u_t(G_t)\), we may suppose that \(H\) is flat over \(S\)
and that \(u_t\) is an isomorphism.

Let \(G^0\) be the connected component of \(G\). By \SGARefLink{sga3:xvi:remark:1:2}{1.2} a), the
restriction \(u|_{G^0}\) is an immersion; denote by \(H^0\) the group
subprescheme of \(H\) that is the image of \(G^0\). Thus
\[
  H_s^0=(H_s)_{\mathrm{red}}^0.
\]

To verify the valuative criterion, we must show that every section
\(h\) of \(H\) over \(S\) such that \(h(S)\) is contained in
\(u(G)\), in the case of an immersion (respectively, every section
\(h\) of \(H\) over \(S\), in the case of a closed immersion), is
the image of a section \(g\) of \(G\) over \(S\).

Let \(C\) be a Cartan subgroup of \(G\), and let
\(D=u(C)\) be the Cartan subgroup of \(H^0\) that is the image of
\(C\). The \(S\)-group
\[
  D'=\operatorname{int}(h)D
\]
has connected fibers, so its underlying space is contained in that
of \(H^0\); on the other hand, since it is smooth over \(S\), it is
reduced, and hence \(D'\) is a smooth group subprescheme of \(H^0\).
The fibers of \(D'\) are Cartan subgroups of the fibers of \(H^0\);
therefore \(D'\) is a Cartan subgroup of \(H^0\), and
\[
  C'=u^{-1}(D')
\]
is a Cartan subgroup of \(G\). But
\(\operatorname{Transpstr}_G(C,C')\) is smooth and surjective over
\(S\) (\SGARefLink{sga3:xii:7.1}{Exposé XII, 7.1} b)). Since \(S\) is henselian with
algebraically closed residue field, there exists \(g\in G(S)\) such
that \(\operatorname{int}(g)C=C'\). Replacing \(h\) by
\(u(g)^{-1}h\), we may suppose that \(h(t)\) normalizes \(D_t\)
(and that \(h(s)\) is the image of an element of the normalizer of
\(C_s\) in \(G_s\), in the case of an immersion). Let
\[
  N=\Norm_G(C).
\]
By hypothesis, \(u|_N\) is an immersion (respectively, a closed
immersion); hence \(h\) is indeed the image of a section of \(N\)
over \(S\), which completes the proof.
\hfill\(\square\)
